\clearpage
\section{Discussion}
\label{sec:discussion}

The experiments support a simple but often neglected claim: unsupervised subtype NMI of shared--private multi-omics models is not a certificate that~$S$ and~$P$ do what their names suggest. This chapter interprets the BRCA multi-seed evidence, relates it to prior work, and discusses limitations and implications.

\subsection{High NMI without role fidelity}

MOCSS~\citep{chen2023mocss} remains a strong clustering system on BRCA. Its failure mode is therefore scientifically informative rather than trivial. Because both shared and private decoders reconstruct the full input under only weak elementwise orthogonality, private codes are free to store subtype-discriminative biology. The empirical signature---$\mathrm{concat}(P)\approx\mathrm{concat}(P,S)>\mathrm{mean}(S)$ with $\Delta\mathrm{X}\,R^{2}\approx0.04$ and subtype-structured private UMAP---matches that architectural incentive. Related representation-learning critiques have long warned that reconstruction objectives identify compressible features, not necessarily the factors specified in a modelling narrative~\citep{bengio2013representation,martens2024disentangling}. Multi-omics shared--private models are not exempt.

The practical consequence for readers of MOCSS-style papers is interpretive caution. A table that reports concatenation NMI near~$0.46$ answers a clustering question; it does not answer whether the shared branch is the carrier of cross-omics subtype biology. In the present multi-seed audit, private codes are at least as subtype-informative as shared codes, and own-view $X\leftarrow P$ exceeds $X\leftarrow S$. That pattern is exactly what dual solo reconstruction encourages.

OMIDIENT~\citep{Safinianaini_Välimäki_Bresson_Gorbonos_Rajamäki_Aaltonen_Marttinen_2025} illustrates a complementary evaluation lesson rather than a dismissal of the method. OMIDIENT's contribution is a stronger private-branch model; its primary purpose was not to certify a new shared definition. Near-zero private subtype NMI and huge $\Delta\mathrm{X}\,R^{2}$ look like perfect separation until own-view $X\leftarrow P$ is examined. Negative reconstructiveness indicates collapse of the Dirichlet private path under the probes of this thesis, not a verified theory of shared content. Hybridising Method~2's joint decoder with OMIDIENT private codes does not remove the collapse signature, which suggests that the Dirichlet private geometry and its KL pressure---not only MOCSS dual reconstruction---can drive uninformative~$P$. The result also cautions against over-trusting single-seed leaderboards: favourable shared NMI in one run can shrink under seeds $\{21,\ldots,25\}$. Technical reproducibility studies in multi-omics make the same methodological point in broader terms~\citep{baiao2025technical,rappoport2021evaluation}.

\subsection{Why minimising private NMI is the wrong objective}

A natural but misleading response to MOCSS-style private leakage is to drive $\mathrm{concat}(P)$ subtype NMI toward zero. Method~1 and Method~1+VICReg show why that proxy fails. Both reduce private NMI relative to MOCSS, and the VICReg ablation even restores a pleasant $S>P$ NMI ordering, yet $\Delta\mathrm{X}\,R^{2}$ stays near zero or negative and concatenation NMI drops. Private codes can be made non-discriminative by collapse, by over-regularisation, or by moving information into brittle shared features that do not integrate well.

Conversely, Method~3's high private NMI is compatible with the best AE-private $\Delta\mathrm{X}\,R^{2}$ when UMAP layouts differ across branches. Assay-specific correlates of subtype are biologically real~\citep{levine2013integrated,parker2009pam50}; a private code that captures them without duplicating shared cross-view geometry is not a bug. Classical integrative factorisations already allowed individual structure to be scientifically meaningful~\citep{lock2013jive,argelaguet2018multi}. Deep models should inherit that nuance rather than treating any private subtype signal as contamination.

The decision rule used in this thesis---jointly read $\Delta\mathrm{X}\,R^{2}$, UMAP role separation, and integration NMI---is therefore not a cosmetic preference. It is required to distinguish complementary private biology from stolen shared signal and from collapsed privacy. Under that rule, Method~1+VICReg's attractive NMI split is insufficient, Method~2 is preferable among simple AE-private fixes, and Method~3 is the clearest overall win on role separation without surrendering competitive concatenation NMI.

\subsection{What the proposed methods achieve}

Method~1 operationalises shared content as cross-view predictability and private content as residual sufficiency. The idea is aligned with classical shared--individual decompositions such as JIVE and MOFA~\citep{lock2013jive,argelaguet2018multi}, translated into deep predictors $f$, $g$, $k$, and $b$. In the BRCA AE setting studied here, the adversarial estimator proved brittle: seed-to-seed variance in $\Delta\mathrm{X}\,R^{2}$ is large, and loss-best configurations under-integrate. The negative result is still useful. It shows that writing an information-theoretic story is easier than estimating it stably with fifteen networks and two-step minimax training.

Method~2 offers a more pragmatic fix: remove dual solo reconstruction, impose VICReg independence~\citep{bardes2022vicreg}, and keep contrastive alignment~\citep{chen2020simple}. Empirically it is the preferred AE-private alternative to the VICReg ablation of Method~1 when both probes and integration matter. The UMAP evidence that shared and private geometries are both structured but differently organised supports the complementary-biology reading of moderate private NMI. Method~2 also has a practical advantage for future extensions: its single-optimiser, all-positive objective is easier to hybridise than Method~1's adversaries, as the OMIDIENT-private companion experiment demonstrates even though that companion inherits private collapse. Nevertheless, Method~2's $\Delta\mathrm{X}\,R^{2}$ ($\approx0.05$) remains close to MOCSS and far below Method~3 ($\approx0.25$). That is a genuine limitation of the model, not a measurement artefact: VICReg only soft-constrains linear $P$--$S$ dependence, and joint reconstruction still permits overlapping content whenever the concatenated decoder succeeds. Without an explicit cross-view definition of sharedness, Method~2 improves hygiene relative to dual solo reconstruction but does not fully enforce the shared--private distinction sought in this thesis.

Method~3 goes further by giving shared codes an explicit spectral geometry target derived from fused patient neighbourhoods~\citep{wang2014snf,ng2002spectral,vonluxburg2007tutorial}. That choice reconnects deep shared--private learning to the classical insight that cross-view agreement lives in sample geometry, not only in per-view feature reconstruction. The large $\Delta\mathrm{X}\,R^{2}$ gain ($\approx0.25$) with competitive concatenation NMI ($\approx0.46$) is the strongest positive result of the thesis. Residual decoding with stop-gradient on~$S$ prevents the shared code from being rewritten as a second full autoencoder, while GRL leakage control and view heads provide additional role regularisation (Appendix~\ref{app:method3}). The comparatively low shared CKA of Method~3 ($\approx0.55$) is consistent with demoting contrastive alignment ($\lambda_{\mathrm{ctr}}=0$ on the loss-best run): spectral regression aligns $\bar{S}$ to~$U$ without forcing the same pairwise linear homogeneity that InfoNCE induces among $S_m$. The spectral ceiling $\mathrm{NMI}(U)\approx0.387$ clarifies a limitation and a research opportunity: Method~3 cannot conjure subtype signal absent from the fused graph, so improving graph construction, affinity kernels, or modality weighting is a direct avenue for further gains.

\subsection{Relation to evaluation standards}

Recent multi-omics benchmarking emphasises robustness, multiple metrics, and careful preprocessing disclosure~\citep{cantini2021benchmarking,rappoport2021evaluation,baiao2025technical,mohr2024navigating}. This thesis adds a missing axis for the shared--private subclass of models: role probes. Reporting only $\mathrm{mean}(S)$ NMI can hide private leakage; reporting only $\mathrm{concat}(P)$ NMI can reward collapse. Cross-view $\Delta\mathrm{X}\,R^{2}$ and own-view $X\leftarrow P$ together catch both failure modes in the BRCA results. Embedding visualisations remain indispensable qualitative checks, especially when numerical gaps are modest, as for Method~2 relative to MOCSS.

Label-free model selection deserves a brief remark. Ranking by core validation loss avoids leaking PAM50 into hyperparameter choice, which is appropriate for an unsupervised claim. It does not optimise $\Delta\mathrm{X}\,R^{2}$ directly. Some sweep cells may therefore dominate the probe metrics while losing on loss rank. The thesis treats that mismatch as an acceptable cost of honest unsupervised selection; future work could explore probe-inspired but still label-free regularisers.

\subsection{Limitations}

Several limitations bound the present conclusions. First, headline results are BRCA-only; KIRC remains to be reported in a later revision, so claims about disease transfer are provisional. Second, PAM50 labels, while standard, are transcriptomic in origin and imperfect as multi-omics ground truth~\citep{parker2009pam50,cantini2021benchmarking}. Third, VICReg enforces linear independence, not full statistical disentanglement~\citep{martens2024disentangling,zbontar2021barlow}. Fourth, Method~1's two-step adversaries may admit stronger estimators, occupancy constraints, or single-step gradient-reversal formulations not explored here. Fifth, figure panels exported from the full repository visualisation may include rows for methods outside this thesis's scope; the narrative intentionally ignores those rows. Sixth, encoder depth and width were fixed to the MOCSS-aligned MLP backbone; capacity ablations might change absolute NMI even if relative role failures persist. Finally, biological interpretation of private complementary signal remains correlative: the thesis does not map private dimensions to specific pathways or CpG modules.

\subsection{Implications for multi-omics representation learning}

For practitioners, the implication is evaluative: published shared--private models should report cross-view reconstruction gaps and own-view private reconstructiveness alongside clustering tables. A minimal reporting set suggested by this work is (i)~NMI/ARI on $\mathrm{mean}(S)$, $\mathrm{concat}(P)$, and $\mathrm{concat}(P,S)$; (ii)~$\Delta\mathrm{X}\,R^{2}$; (iii)~own-view $X\leftarrow P$; and (iv)~at least one embedding figure coloured by subtype for both branches.

For method designers, the implication is architectural. If~$P$ is meant to be residual, do not train it as a second full autoencoder of~$X$. If~$S$ is meant to be shared, give it an explicit cross-view target---predictability, contrastive agreement, or spectral geometry---that reconstruction alone does not provide. Classical factorisations already embodied these principles in linear form~\citep{lock2013jive,shen2009integrative,argelaguet2018multi}; deep models must re-enforce them rather than assume that naming latents ``shared'' and ``private'' is sufficient. Among the concrete designs studied here, Method~3 currently offers the best empirical trade-off on BRCA, with Method~2 as the simpler AE-private baseline recommended when spectral precomputation is undesirable.
